\section{Related Work}\label{sec:related}

\paragraph{Equivariant GNNs for atomistic modelling.}
Invariant architectures such as SchNet~\citep{schutt2017schnet} and DimeNet++~\citep{gasteiger2020dimenetpp} encode atomic environments through radial and angular features but discard directional information at the representation level. Equivariant models restore this by constructing features that transform as irreducible representations of $\mathrm{SO}(3)$: NequIP~\citep{batzner2022nequip} propagates tensor field features through learned Clebsch--Gordan tensor products, MACE~\citep{batatia2022mace} achieves data efficiency via higher-order symmetric contractions, and PaiNN~\citep{schutt2021painn}, TorchMD-NET~\citep{tholke2022torchmdnet}, and ViSNet~\citep{wang2024visnet} operate on equivariant vector features ($L{=}1$) with favourable speed but limited angular resolution. EquiformerV2~\citep{liao2024equiformerv2} combines equivariant attention with eSCN convolutions~\citep{passaro2023escn} supporting tensor products up to $L{=}6$ at substantial computational cost. FAENet~\citep{duval2023faenet} achieves equivariance through frame averaging without explicit group representations, excelling at scalar prediction but unable to predict forces natively. All CG-based methods share a fundamental bottleneck: tensor product cost grows combinatorially with $L$.

\paragraph{CG-free equivariant architectures.}
Several approaches achieve equivariant tensor interactions without Clebsch--Gordan coefficients. ICTP~\citep{zaverkin2024ictp} constructs products of irreducible Cartesian tensors (symmetric traceless tensors) within a MACE-like architecture, reporting results up to $L{=}3$ on rMD17, MD22, 3BPA, and acetylacetone. GotenNet~\citep{aykent2025gotennet} uses geometry-aware tensor attention via inner products on high-degree steerable features. Gaunt tensor products~\citep{luo2024gaunt} exploit spherical-harmonic addition theorems to reduce CG complexity from $\mathcal{O}(L^6)$ to $\mathcal{O}(L^3)$. TensorNet~\citep{simeon2023tensornet} replaces CG products with rank-2 Cartesian tensor operations, CACE~\citep{cheng2024cace} reformulates the atomic cluster expansion in Cartesian coordinates, and HotPP~\citep{wang2024n} passes reducible Cartesian tensor messages of arbitrary rank. The mathematical foundation underlying this entire line of work is the classical equivalence between symmetric traceless Cartesian tensors and $\mathrm{SO}(3)$ spherical irreps \citep{coope1965irreducible,applequist1989traceless,thorne1980multipole,stone2013theory}, recently placed on rigorous algorithmic footing via an orthonormal change-of-basis construction up to rank 9~\citep{shao2025orthonormal}. Our Theorem~\ref{thm:cg_completeness} operationalises this equivalence at the architectural level through the specific combination of the $\mathrm{Cl}(3,0)$ geometric product with closed-form STF operations, enabling a CG-free forward pass that invokes neither CG coefficients nor Wigner-$D$ matrices nor spherical-harmonic evaluation. CliffordSTF shares the STF primitive with ICTP but differs in a substantive way: ICTP replaces CG products with a single Cartesian substrate, while CliffordSTF runs the Clifford multivector and STF tracks in parallel and couples them through a bidirectional cross-track bilinear. This two-substrate design is architecturally distinct, and our ablation shows the cross-track coupling is essential for end-to-end training stability, with its removal approximately doubling both the mean and standard deviation of energy MAE across seeds (Supplementary~\ref{supp:ictp}).

\paragraph{Clifford algebra, and force directional fidelity.}
Clifford algebra has been applied to equivariant networks across several domains: geometric Clifford algebra networks~\citep{ruhe2023gcan} introduced $\mathrm{Pin}(p,q,r)$ group layers, Clifford group equivariant neural networks~\citep{ruhe2023cgenn} extended equivariance to the full Clifford group, Clifford neural layers~\citep{brandstetter2022clifford} applied multivector feature maps to PDE surrogates, GATr~\citep{brehmer2023geometric} built a transformer over the projective algebra $\mathrm{Cl}(3,0,1)$, Clifford-steerable CNNs~\citep{zhdanov2024clifford} construct equivariant kernels on pseudo-Riemannian manifolds, and CG-EGNN~\citep{tran2024cgegnn} combined Clifford layers with high-order message passing. No prior work identifies that the geometric product of two vectors in $\mathrm{Cl}(3,0)$ realises only the $L{=}0$ and $L{=}1$ components of $\mathbf{1}\otimes\mathbf{1} = \mathbf{0}\oplus\mathbf{1}\oplus\mathbf{2}$, leaving $L \geq 2$ inaccessible through the per-edge bilinear regardless of depth. Higher algebraic \emph{grade} in $\mathrm{Cl}(3,0)$ does not correspond to higher angular momentum $L$: grade-2 bivectors transform as $L{=}1$ pseudovectors under $\mathrm{SO}(3)$, not $L{=}2$. The same deficiency applies to the projective algebra $\mathrm{Cl}(3,0,1)$ used in GATr, where the additional null basis encodes translational rather than angular content and the $\mathrm{SO}(3)$ decomposition remains $2\cdot L{=}0 \oplus 2\cdot L{=}1$ (Supplementary~\ref{supp:gatr}). Standard magnitude metrics can obscure these representational differences: low force MAE does not reliably predict simulation quality~\citep{fu2023forces}, and hyperparameters optimising force accuracy can degrade energy accuracy~\citep{kovacs2023mace}. Force cosine similarity, part of the OC20 benchmark since its introduction but not standardly reported on rMD17, exposes the gap directly; we bring it to small-molecule evaluation and observe that $L \leq 1$ direct-force architectures cluster at poor directional alignment while $L \geq 2$ models achieve substantially better values, consistent with our theoretical framing.